\chapter{MÔ HÌNH HỆ THỐNG VÀ BÀI TOÁN TỐI ƯU}
\label{chap:system-model}

\section{Kiến trúc hệ thống}
Luận văn xét đường xuống của một tế bào đơn gồm một trạm gốc (Base Station, BS) SISO, một STAR--RIS thụ động có $N$ phần tử và $K=4$ người dùng SISO. Người dùng 1 và 3 nằm ở miền phản xạ; người dùng 2 và 4 nằm ở miền truyền qua. Việc gán phía được cố định trong toàn bộ kịch bản để tập trung đánh giá bài toán phân bổ tài nguyên.

BS phát một luồng chung và bốn luồng riêng theo RSMA. Mỗi người dùng nhận tín hiệu từ đường trực tiếp BS--người dùng và đường ghép tầng BS--STAR--RIS--người dùng. STAR--RIS vận hành theo ES, vì vậy mỗi phần tử đồng thời tạo thành phần truyền qua và phản xạ.

\ThesisFigure{figures/pdf/chapter2_fig01_system_model_3d.pdf}
{Mô hình hệ thống SISO STAR--RIS hỗ trợ RSMA}
{fig:ch2-system-model}[0.95\linewidth]
\figuresource{Tác giả xây dựng từ mô hình vật lý và cấu hình thực nghiệm của đề tài (2026)}

Hình~\ref{fig:ch2-system-model} thể hiện STAR--RIS dưới dạng tấm nghiêng theo phối cảnh. Miền phản xạ nằm bên trái, miền truyền qua nằm bên phải. Đường xám nét đứt biểu diễn kênh trực tiếp; đường xanh dương từ BS đến STAR--RIS cùng các đường xanh lá hoặc cam tạo thành kênh ghép tầng. Hình cũng cho thấy một luồng chung $s_c$ và bốn luồng riêng $s_1,\ldots,s_4$ được truyền đồng thời.

\section{Các giả định nghiên cứu}
Các giả định được sử dụng nhất quán trong mô hình và mô phỏng gồm:
\begin{itemize}
  \item CSI hoàn hảo được cung cấp cho bộ ra quyết định;
  \item các hệ số kênh được sinh độc lập theo phân bố Gaussian phức;
  \item STAR--RIS thụ động, vận hành theo ES và không khuếch đại tín hiệu;
  \item pha truyền và pha phản xạ được điều khiển độc lập;
  \item SIC lý tưởng, không có nhiễu dư sau khi loại bỏ luồng chung;
